\uuid{D8uh}
\titre{Étude de suite}
\auteur{bodin ; hanani ; mzari}
\organisation{exo7}
\licence{CC BY-NC-SA 4.0}
\datecreate{2020-01-20}
\niveau{L1}
\module{Analyse}
\chapitre{Suite}
\sousChapitre{Autre}
\difficulte{3}
\theme{Suites, étude de suite}
\exo7qcmid{lille-1.suites.029}

\contenu{
\begin{question}
\qclassification{121.00}
\qtags{motcle=étude de suite}
Soit $(u_n)$ une suite croissante. On suppose que $\displaystyle u_n+\frac{1}{n+1}\leq u_{n+1}$ pour tout $n\geq 0$. Quelles sont les bonnes réponses ?
\begin{answers}
    \bad{$(u_n)$ est majorée.}
    \good{$(u_n)$ est divergente.}
    \bad{$(u_n)$ est convergente et $\displaystyle \lim _{n\to +\infty}u_n\geq 0$.}
    \bad{$u_n=0$ pour tout $n\geq 1$.}
\end{answers}
\begin{explanations}
On vérifie par récurrence que, pour tout $n\geq 1$,
$$\displaystyle u_0+1+\frac{1}{2}+\dots +\frac{1}{n}\leq u_n.$$
Donc, $(u_n)$ n'est pas majorée car sinon, il en serait de même pour la suite de terme général $\displaystyle v_n=1+\frac{1}{2}+\dots +\frac{1}{n}$ et on sait que $\displaystyle \lim _{n\to +\infty}v_n=+\infty$.
\end{explanations}
\end{question}
}
